\chapter{Összegzés}
\label{ch:sum}
A dolgozatban egy nyelvtanulást segítő webalkalmazás fejlesztését ismertettem, amely japán és angol szövegek feldolgozására, szókincs- és kanji tanulásra, valamint a folyamat követésére szolgál. A kiinduló problémát az jelentette, hogy a hagyományos szótanuló eszközök nem feltétlenül képesek hatékonyan támogatni a szövegközpontú tanulást, megfelelően kezelni a japán nyelv sajátosságait, és nem veszik kellőképpen figyelembe a tanulók egyéni előrehaladását. Célom egy olyan rendszer létrehozása volt, amely automatikusan feldolgozza a szövegeket, azonosítja a bennük szereplő nyelvi elemeket és ezekből tanulókártyákat generál, miközben a tanulási statisztikákat is folyamatosan rögzíti.

A megvalósításhoz backendként a Django keretrendszert használtam, amely az adatmodellek, nézetek, API-k és az adminisztrációs felület kezelését biztosítja. A szövegfeldolgozásért egy külön \texttt{processor} modul felel, amely a japán és angol tokenizálást, mondatok felbontását, nyelvfelismerést és validációt végzi. A \texttt{dashboard} modul adja a felhasználói felületet, ahol a tanuló szövegeket, szavakat és kanji karaktereket böngészhet, statisztikákat tekinthet meg, illetve saját szótárakat hozhat létre. A \texttt{dictionary} modulban valósítottam meg az Anki alkalmazáshoz hasonló tanulási mechanizmust, ahol a szavak és kanji karakterek ismétlési intervallumai, státuszai, valamint a napi tanulási mutatók a felhasználói válaszok alapján automatikusan frissülnek.

A projekt továbbfejlesztésére számos lehetőség kínálkozik. A jelenlegi SQLite adatbázis jelentősen nagyobb mennyiségű adat kezelésére lecserélhető PostgreSQL-re, ami nagyobb teljesítményt és fejlettebb adatkezelési funkciókat biztosítana. A végleges környezetben fejlesztői szerver helyett célszerűbb biztonságosabb és stabilabb üzemeltetést lehetővé tevő, egyszerűen telepíthető megoldást választani. A rendszer nyelvi moduljai bővíthetők további nyelvekkel, ehhez új tokenizáló és statisztikai algoritmusok integrálása szükséges. A felhasználókezelés fejlesztésével (regisztráció, jogosultságok, profilok) lehetőség nyílna több tanuló egyidejű kezelésére, valamint személyre szabott tanulási statisztikákra és közösségi funkciókra is.

A projekt során sikerült létrehozni egy olyan nyelvtanuló alkalmazást, amely rugalmasan bővíthető, könnyen testre szabható, és valóban segíti a szövegalapú tanulást, legyen szó japánról vagy angolról. A fejlesztés közben azt tapasztaltam, hogy a modern webes technológiák és a nyelvi feldolgozó algoritmusok jól kiegészítik egymást, és tényleg képesek támogatni a tanulók egyéni fejlődését.